\subsection{Distribuição radial uniforme dos terminais}
\label{subsec:distribuicao-radial-uniforme}

Na Bateria Test 3, a posição inicial de cada equipamento de usuário
(\textit{User Equipment}--UE) é sorteada em coordenadas polares. O ângulo
$\Theta$ é uniforme no intervalo $[0,2\pi)$ e a distância radial $R$ é uniforme
entre $r_{\min}=10\,\mathrm{m}$ e $r_{\max}=500\,\mathrm{m}$:
\begin{equation}
  \Theta \sim \mathcal{U}(0,2\pi),
  \qquad
  R \sim \mathcal{U}(r_{\min},r_{\max}).
\end{equation}
As coordenadas cartesianas são obtidas por
\begin{equation}
  X=R\cos\Theta,
  \qquad
  Y=R\sin\Theta,
  \qquad
  Z=1{,}5\,\mathrm{m}.
\end{equation}
Consequentemente, a função densidade da distância é
\begin{equation}
  f_R(r)=\frac{1}{r_{\max}-r_{\min}},
  \quad r_{\min}\leq r\leq r_{\max}.
\end{equation}
Esse desenho atribui a mesma probabilidade a intervalos radiais de igual
largura. Por exemplo, desconsiderando efeitos de borda, um anel de $10\,$m
próximo à estação base contém a mesma probabilidade que um anel de $10\,$m
próximo ao limite de cobertura. Como a área do anel cresce com o raio, a
densidade de UEs por unidade de área não é constante. A densidade espacial
resultante é proporcional a
\begin{equation}
  f_{X,Y}(x,y)=
  \frac{1}{2\pi(r_{\max}-r_{\min})\sqrt{x^2+y^2}},
\end{equation}
no anel $r_{\min}\leq\sqrt{x^2+y^2}\leq r_{\max}$.

Essa escolha é intencional: o objetivo não é reproduzir uma população
uniforme por metro quadrado, mas obter cobertura aproximadamente equilibrada
das diferentes faixas de distância e, assim, submeter os escalonadores a uma
amostra ampla de condições de canal, do centro à borda. Em cada
\texttt{rngRun}, uma nova topologia é sorteada. Dentro da mesma
\texttt{rngRun}, RR, PF, MR e QoS recebem exatamente as mesmas coordenadas,
permitindo contrastes pareados. Os UEs permanecem estáticos durante os
$30\,$s simulados; portanto, cada replicação representa um instantâneo de
implantação, e não uma trajetória de mobilidade.

Uma distribuição uniforme em área exigiria, em contraste,
\begin{equation}
 R=\sqrt{r_{\min}^{2}+U(r_{\max}^{2}-r_{\min}^{2})},
 \qquad U\sim\mathcal{U}(0,1),
\end{equation}
e produziria mais terminais nos anéis externos. Essa alternativa não é usada
na Bateria Test 3 e deve ser tratada como outro cenário experimental.
